% Luc Destombe --- licence CC BY-NC-SA 4.0
% https://creativecommons.org/licenses/by-nc-sa/4.0/deed.fr
% Fichier autonome : compiler avec pdflatex, deux fois (signets, nombre de pages).
\def\lyceeversion{1.0}
\documentclass[a4paper,10pt]{article}
\makeatletter
\newif\iflycee@niveaux
\lycee@niveauxtrue
\newif\iflycee@palatino
\lycee@palatinotrue

\RequirePackage[utf8]{inputenc}
\RequirePackage[T1]{fontenc}
\RequirePackage[french]{babel}
\RequirePackage{etoolbox}

\newcommand{\ShorthandsOff}{\@ifundefined{shorthandoff}{}{\shorthandoff{;:!?}}}
\newcommand{\ShorthandsOn}{\@ifundefined{shorthandon}{}{\shorthandon{;:!?}}}
\ShorthandsOff
\AtBeginDocument{\ShorthandsOn}
\AtBeginEnvironment{tikzpicture}{\ShorthandsOff}

\RequirePackage{geometry}
\RequirePackage{fancyhdr}
\RequirePackage{lastpage}
\RequirePackage{multicol}
\RequirePackage{array}
\RequirePackage{enumitem}
\RequirePackage{xcolor}
\RequirePackage{needspace}

\RequirePackage{amsmath,amssymb}
\RequirePackage{mathpazo}
\renewcommand{\ttdefault}{lmtt}
\DeclareMathSizes{10.95}{10.95}{8}{5.5}
\begingroup\catcode`\_=\active \catcode`\^=\active
\gdef_#1{\sb{\scriptscriptstyle #1}}
\gdef^#1{\sp{\scriptscriptstyle\lycee@moinsepais #1}}
\endgroup
\newcommand\lycee@moins{\mathbin{%
  \setbox\z@\hbox{$\m@th\scriptscriptstyle\lycee@moinsorig$}%
  \dimen@=.55\wd\z@ \advance\dimen@-.5pt
  \hbox to.75\wd\z@{\hss$\m@th\scriptscriptstyle\vcenter{\hbox{%
    \tikz\draw[line width=.5pt,line cap=round](0,0)--(\the\dimen@,0);}}$\hss}}}
\begingroup\lccode`\~=`\- \lowercase{\endgroup\let~}\lycee@moins
\newcommand\lycee@moinsepais{\mathcode`\-="8000 }
\AtBeginDocument{\mathchardef\lycee@moinsorig=\mathcode`\-\relax}
\AtBeginDocument{\catcode`\_=12 \mathcode`\_="8000
  \catcode`\^=12 \mathcode`\^="8000 }
\newcommand\lycee@exposants{%
  \fontdimen13\textfont2=.48\fontdimen6\textfont2
  \fontdimen14\textfont2=.48\fontdimen6\textfont2
  \fontdimen15\textfont2=.48\fontdimen6\textfont2
  \fontdimen13\scriptfont2=.48\fontdimen6\scriptfont2
  \fontdimen14\scriptfont2=.48\fontdimen6\scriptfont2
  \fontdimen15\scriptfont2=.48\fontdimen6\scriptfont2}
\every@math@size\expandafter{\the\every@math@size\lycee@exposants}
\newlength{\retraitcalculs}
\setlength{\retraitcalculs}{2em}

\RequirePackage{tikz}
\usetikzlibrary{arrows.meta,calc,babel,shapes.misc}
\RequirePackage[most]{tcolorbox}
\tcbuselibrary{skins,breakable}

\definecolor{coulDef}{HTML}{1F4E79}
\definecolor{coulProp}{HTML}{7030A0}
\definecolor{coulRem}{HTML}{C55A11}
\definecolor{coulEx}{HTML}{2E7D32}
\definecolor{coulExo}{HTML}{B03A2E}
\definecolor{coulPart}{HTML}{1F4E79}
\definecolor{coulMeth}{HTML}{1B6E4F}
\definecolor{coulCourbe}{HTML}{0B5ED7}
\definecolor{coulCourbeB}{HTML}{C00000}
\definecolor{coulCourbeC}{HTML}{2E7D32}
\definecolor{coulGrille}{HTML}{8C8C8C}
\definecolor{coulRep}{HTML}{1B6E4F}
\definecolor{coulBar}{HTML}{2E7D32}

\newcommand{\R}{\mathbb{R}}
\newcommand{\Rs}{\mathbb{R}^{*}}
\renewcommand{\le}{\leqslant}
\renewcommand{\ge}{\geqslant}
\renewcommand{\approx}{\simeq}
\newcommand{\vect}[1]{\overrightarrow{#1}}
\newcommand{\norme}[1]{\left\lVert #1 \right\rVert}
\newcommand{\intervalle}[2]{\left[#1\,;#2\right]}
\RequirePackage{cancel}
\renewcommand{\CancelColor}{\color{coulExo}}
\newcommand{\fois}[1]{\times{\color{coulExo}#1}}
\newcommand{\barre}[1]{\cancel{#1}}

\tikzset{grille/.style={coulGrille,line width=0.4pt}}
\newcommand{\quadrillage}[4]{\draw[grille,step=1] (#1,#2) grid (#3,#4);}
\tikzset{
  croix/.style={cross out,draw,line width=0.6pt,inner sep=0pt,minimum size=4pt},
  croixdroite/.style={croix,rotate=45,line width=0.8pt},
}
\newcommand{\pt}[3]{\path (#1) node[croix]{} node[#2,font=\small,inner sep=2.5pt]{$#3$};}

\newcommand{\lycee@repere}[4]{%
  \draw[grille,step=1] (#1,#2) grid (#3,#4);
  \draw[-{Stealth[length=2mm]},thick] (#1,0)--({#3+0.4},0) node[below left=-1pt]{$x$};
  \draw[-{Stealth[length=2mm]},thick] (0,#2)--(0,{#4+0.4}) node[below left=-1pt]{$y$};
  \node[below left,font=\scriptsize,inner sep=1.5pt] at (0,0) {$0$};}
\newcommand{\lycee@gradx}[1]{\foreach \k/\l in {#1}{%
  \draw (\k,0.07)--(\k,-0.07) node[below,font=\scriptsize,inner sep=1.5pt]{$\l$};}}
\newcommand{\lycee@grady}[1]{\foreach \k/\l in {#1}{%
  \draw (0.07,\k)--(-0.07,\k) node[left,font=\scriptsize,inner sep=1.5pt]{$\l$};}}
\newcommand{\lycee@intff}[2]{\left[#1\,;#2\right]}
\newcommand{\lycee@intfo}[2]{\left[#1\,;#2\right[}
\newcommand{\lycee@intof}[2]{\left]#1\,;#2\right]}
\newcommand{\lycee@intoo}[2]{\left]#1\,;#2\right[}
\newcommand{\lycee@ens}[1]{\left\{#1\right\}}
\AtBeginDocument{%
  \providecommand{\repere}{\lycee@repere}%
  \providecommand{\gradx}{\lycee@gradx}%
  \providecommand{\grady}{\lycee@grady}%
  \providecommand{\Cf}{\mathcal{C}\sb{\scriptscriptstyle f}}%
  \providecommand{\Cg}{\mathcal{C}\sb{\scriptscriptstyle g}}%
  \providecommand{\intff}{\lycee@intff}%
  \providecommand{\intfo}{\lycee@intfo}%
  \providecommand{\intof}{\lycee@intof}%
  \providecommand{\intoo}{\lycee@intoo}%
  \providecommand{\ens}{\lycee@ens}}

\newlist{questions}{enumerate}{3}
\setlist[questions,1]{label=\arabic*\textdegree),leftmargin=*,itemsep=2pt,topsep=3pt}
\setlist[questions,2]{label=\alph*),leftmargin=*,itemsep=1pt,topsep=2pt}
\setlist[questions,3]{label=\roman*),leftmargin=*,itemsep=1pt,topsep=1pt}
\setlist[itemize]{label=\textbullet,leftmargin=*,itemsep=2pt,topsep=3pt}

\newcommand{\besoinplace}[1]{\par\penalty\@M\begingroup
  \dimen@=#1\relax\dimen@ii=\pagegoal\advance\dimen@ii-\pagetotal
  \ifdim\dimen@>\dimen@ii\ifdim\dimen@ii>\z@\vfil\fi\break\fi\endgroup}

\newcommand{\lycee@pied}{\thepage/\pageref{LastPage}}
\newcommand{\entete}[2]{%
  \pagestyle{fancy}\fancyhf{}%
  \fancyhead[L]{\footnotesize\color{coulPart}\textbf{#1}}%
  \fancyhead[R]{\footnotesize\color{coulPart}\textbf{#2}}%
  \fancyfoot[C]{\footnotesize\lycee@pied}%
  \renewcommand{\headrulewidth}{0.4pt}%
  \renewcommand{\headrule}{\hbox to\headwidth{%
    \color{coulPart!40}\leaders\hrule height \headrulewidth\hfill}}}

\RequirePackage{ccicons}
\newcommand{\lycee@licence}{{\scriptsize\color{gray}%
  \copyright\ Luc Destombe --- \ccbyncsa\ CC BY-NC-SA 4.0}}
\newif\iflycee@licence \lycee@licencetrue
\newcommand{\sanslicence}{\lycee@licencefalse}
\AtBeginDocument{\iflycee@licence\AddToHookNext{shipout/foreground}{%
  \put(0,0){\rlap{\kern\dimexpr1in+\hoffset+\oddsidemargin+\textwidth\relax
    \llap{\raisebox{-\dimexpr1in+\voffset+\topmargin+\headheight+\headsep
      +\textheight+\footskip\relax}{\lycee@licence}}}}}\fi}

\newcommand{\lycee@site}{faire-des-maths.fr}
\newcommand{\lycee@version}{{\scriptsize\color{gray}\lycee@site
  \ifdefined\lyceeversion\ --- Version~\lyceeversion\fi}}
\AtBeginDocument{\iflycee@licence\AddToHookNext{shipout/foreground}{%
  \put(0,0){\rlap{\kern\dimexpr1in+\hoffset+\oddsidemargin\relax
    \raisebox{-\dimexpr1in+\voffset+\topmargin+\headheight+\headsep
      +\textheight+\footskip\relax}{\lycee@version}}}}\fi}

\newcommand{\formulesserrees}{%
  \setlength{\abovedisplayskip}{3pt}\setlength{\belowdisplayskip}{3pt}%
  \setlength{\abovedisplayshortskip}{2pt}\setlength{\belowdisplayshortskip}{3pt}}

\setlength{\parindent}{0pt}

\RequirePackage[implicit=false,hidelinks,pdfencoding=auto,psdextra,bookmarksopen,
  bookmarksopenlevel=1,pdfcreator={},pdfproducer={}]{hyperref}
\RequirePackage{bookmark}
\hypersetup{pdfauthor={Luc Destombe},pdfkeywords={CC BY-NC-SA 4.0}}
\newcounter{lycee@signet}
\newcommand{\signet}[3][]{\stepcounter{lycee@signet}%
  \edef\lycee@dest{\if\relax\detokenize{#1}\relax
    signet.\arabic{lycee@signet}\else#1\fi}%
  \hypertarget{\lycee@dest}{}\bookmark[level=#2,dest=\lycee@dest]{#3}}
\pdfstringdefDisableCommands{%
  \def\R{R}\def\vect#1{#1}\def\Cf{Cf}\def\Cg{Cg}\def\fois#1{×#1}%
  \def\barre#1{#1}\def\intervalle#1#2{[#1;#2]}\def\intff#1#2{[#1;#2]}%
  \def\textsuperscript#1{#1}\def\dfrac#1#2{#1/#2}\def\frac#1#2{#1/#2}%
  \def\vec#1{#1⃗}}%
\RequirePackage[official]{eurosym}

\geometry{top=0.8cm,bottom=1.3cm,left=1cm,right=1cm,footskip=0.6cm}

\pagestyle{fancy}
\fancyhf{}
\renewcommand{\headrulewidth}{0pt}
\fancyfoot[C]{\footnotesize Page \thepage{} / \pageref{LastPage}}

\newcommand{\titrefiche}[2]{%
  \begin{center}
    \Large\textbf{#1}\\[2pt]
    \large\textbf{#2}\\[2pt]
  \end{center}
  \vspace{0.10cm}}

\setlength{\columnseprule}{0.4pt}
\setlength{\columnsep}{15pt}
\renewcommand{\columnseprulecolor}{\color{black!45}}
\setlength{\multicolsep}{2pt}

\newcounter{partiefiche}
\renewcommand{\thepartiefiche}{\Roman{partiefiche}}
\newif\ifapresbandeau
\newcommand{\lycee@bandeau}[1]{%
  \par\penalty-600
  \vspace{0.08cm}%
  \noindent\colorbox{coulPart}{%
    \parbox{\dimexpr\linewidth-2\fboxsep\relax}{%
      \color{white}\bfseries\raggedright\strut#1\strut}}%
  \nopagebreak\par\nopagebreak
  \vspace{0.06cm}\nopagebreak
  \apresbandeautrue}
\newcommand{\partiefiche}[1]{%
  \refstepcounter{partiefiche}%
  \lycee@bandeau{\signet{1}{\thepartiefiche{} --- #1}\thepartiefiche{} --- #1}}
\newcommand{\partiefichesansnum}[1]{\lycee@bandeau{\signet{1}{#1}#1}}

\newcounter{exo}
\newlength{\espaceexo}\setlength{\espaceexo}{7pt}
\newlength{\espacetitre}\setlength{\espacetitre}{2pt}
\newcommand{\exo}[1][\relax]{%
  \par
  \ifapresbandeau\apresbandeaufalse\else\penalty-150\addvspace{\espaceexo}\fi
  \refstepcounter{exo}%
  \noindent\signet[exercice-\arabic{exo}]{2}{Exercice \arabic{exo}}%
  \textbf{\color{coulExo}\underline{Exercice \theexo}}%
  \ifx\relax#1\relax\else\textbf{ : #1}\fi
  \par\nopagebreak\vspace{\espacetitre}\nopagebreak
  \ignorespaces}

\setlist[enumerate]{nosep,leftmargin=*,itemsep=2pt,topsep=2pt}
\setlist[itemize]{nosep,leftmargin=*,topsep=2pt}
\newcommand{\choix}[3]{\par\hspace*{1em}%
  \makebox[0.3\linewidth][l]{a) #1}\makebox[0.3\linewidth][l]{b) #2}c) #3}
\newcommand{\deux}[2]{\par\hspace*{0.6em}%
  \makebox[0.48\linewidth][l]{#1}#2}
\newcommand{\trois}[3]{\par\hspace*{0.6em}%
  \makebox[0.32\linewidth][l]{#1}\makebox[0.32\linewidth][l]{#2}#3}

  \renewcommand{\exo}[1][\relax]{%
    \ifapresbandeau\apresbandeaufalse\else\par\penalty-150\fi
    \refstepcounter{exo}%
    \vspace{0.05cm}%
    \noindent\signet[exercice-\arabic{exo}]{2}{Exercice \arabic{exo}}%
    \textbf{\color{coulExo}\underline{Exercice \theexo}}%
    \ifx\relax#1\relax\else\textbf{ : #1}\fi%
    \par\nopagebreak
    \ignorespaces}
  \newcommand{\rappel}[1]{%
    \par\vspace{0.06cm}\nopagebreak
    \noindent\fcolorbox{black!35}{black!5}{%
      \parbox{\dimexpr\linewidth-2\fboxsep-2\fboxrule\relax}{%
        \small\textbf{Rappel.} #1}}%
    \par\vspace{0.06cm}\nopagebreak}
  \setlist[enumerate]{nosep,leftmargin=*}
  \setlist[itemize]{nosep,leftmargin=*}
  \setlength{\multicolsep}{3pt plus 1pt minus 1pt}
  \setlength{\premulticols}{10pt}
  \setlength{\postmulticols}{10pt}
  \newlist{expr}{enumerate}{1}
  \setlist[expr]{label={\Alph*\ $=$},nosep,leftmargin=*,itemsep=0pt}

\renewenvironment{center}{\par\addvspace{3pt}\centering}{\par\addvspace{3pt}}
\formulesserrees
\AtBeginDocument{\formulesserrees}
\clubpenalty=10000
\widowpenalty=10000
\displaywidowpenalty=10000
\brokenpenalty=10000
\setlength{\parskip}{0pt}

\RequirePackage{ragged2e}
\setlength{\RaggedRightRightskip}{0pt plus 1fil}
\AtBeginDocument{\RaggedRight\binoppenalty=10000 \relpenalty=10000 }
\g@addto@macro\@parboxrestore{\RaggedRight}
\makeatother
\usepackage{tkz-tab}
\geometry{top=0.8cm,bottom=1.3cm,left=1cm,right=1cm,footskip=0.7cm,headheight=12pt}

\newcommand{\e}{\mathrm{e}}

\begin{document}
\titrefiche{1\textsuperscript{re} spécialité --- Fonction exponentielle}{Fiche d'exercices du chapitre}

\begin{multicols}{2}

\partiefiche{Définition, relation fonctionnelle, notation $\e^{x}$}

\exo[Écrire sous la forme $\e^{n}$]
Écrire chaque nombre sous la forme $\e^{n}$, où $n$ est un entier relatif.
\begin{enumerate}[label=\alph*)]
  \item $A=\e^{3}\times\e^{6}$
  \item $B=\e^{2}\times\e^{-5}$
  \item $C=\dfrac{\e^{9}}{\e^{4}}$
  \item $D=\dfrac{\e^{-1}}{\e^{-6}}$
  \item $E=\left(\e^{4}\right)^{3}$
  \item $F=\dfrac{1}{\e^{-7}}$
  \item $G=\left(\e^{-2}\right)^{4}\times\e^{5}$
\end{enumerate}

\exo[Simplifier une expression]
Pour tout nombre réel $x$, écrire chaque expression sous la forme $\e^{\dots}$.
\begin{enumerate}[label=\alph*)]
  \item $A=\e^{2x}\times\e^{x+3}$
  \item $B=\dfrac{\e^{4x-1}}{\e^{x+2}}$
  \item $C=\left(\e^{x}\right)^{3}\times\e^{-2x}$
  \item $D=\dfrac{\e^{x+5}}{\e^{3x-1}}$
  \item $E=\left(\e^{-x}\right)^{2}\times\e^{3x}$
  \item $F=\dfrac{\e^{3x}}{\e^{x}\times\e^{2x-4}}$
\end{enumerate}

\exo[Développer]
Pour tout nombre réel $x$, développer et simplifier chaque expression.
\begin{enumerate}[label=\alph*)]
  \item $A=\left(\e^{x}+1\right)^{2}$
  \item $B=\left(\e^{x}-3\right)^{2}$
  \item $C=\left(\e^{x}+2\right)\left(\e^{x}-2\right)$
  \item $D=\e^{x}\left(\e^{x}+\e^{-x}\right)$
  \item $E=\left(\e^{x}+\e^{-x}\right)^{2}$
\end{enumerate}

\exo[Vrai ou faux]
Pour chaque égalité, dire si elle est vraie pour tout nombre réel $x$, et le justifier
par une phrase qui cite la règle utilisée (ou un contre-exemple).
\begin{enumerate}[label=\alph*)]
  \item $\e^{2}\times\e^{3}=\e^{6}$
  \item $\e^{-x}=-\e^{x}$
  \item $\left(\e^{x}\right)^{2}=\e^{2x}$
  \item $\e^{x}+\e^{x}=\e^{2x}$
  \item $\e^{0}=0$
  \item $\e^{x}\times\e^{-x}=1$
\end{enumerate}

\exo[Une expression constante]
Pour tout nombre réel $x$, on pose
\[
  f(x)=\left(\e^{x}+\e^{-x}\right)^{2}-\left(\e^{x}-\e^{-x}\right)^{2}.
\]
\begin{enumerate}
  \item Calculer $f(0)$ et $f(1)$ à la calculatrice.
  \item Développer $\left(\e^{x}+\e^{-x}\right)^{2}$ et
        $\left(\e^{x}-\e^{-x}\right)^{2}$.
  \item En déduire que $f(x)$ ne dépend pas de $x$.
\end{enumerate}

\partiefiche{Équations et inéquations avec des exponentielles}

\exo[Résoudre une équation]
Résoudre dans $\R$ chaque équation.
\begin{enumerate}[label=\alph*)]
  \item $\e^{3x+1}=\e^{x+7}$
  \item $\e^{5x-10}=1$
  \item $\e^{2x+3}=\e^{1-x}$
  \item $\e^{x}=\e$
  \item $\e^{4-x}=-2$
\end{enumerate}

\exo[Équations et second degré]
Résoudre dans $\R$ chaque équation.
\begin{enumerate}[label=\alph*)]
  \item $\e^{x^{2}}=\e^{2x+8}$
  \item $\e^{x^{2}-5x}=1$
  \item $\e^{x^{2}+1}=\e^{x}$
  \item $\e^{2x}\times\e^{x^{2}}=\e^{3}$
\end{enumerate}

\exo[Résoudre une inéquation]
Résoudre dans $\R$ chaque inéquation.
\begin{enumerate}[label=\alph*)]
  \item $\e^{2x+5}\le\e^{x+8}$
  \item $\e^{3-x}>1$
  \item $\e^{4x}\ge\e^{x-6}$
  \item $\e^{x^{2}}<\e^{x+2}$
  \item $\e^{x^{2}}\ge\e^{4}$
\end{enumerate}

\exo[Produit par une exponentielle]
Une exponentielle n'est jamais nulle et toujours strictement positive. En s'en servant,
résoudre dans $\R$ :
\begin{enumerate}[label=\alph*)]
  \item $(x-3)\,\e^{x}=0$
  \item $(2x+1)\,\e^{x}>0$
  \item $\left(x^{2}-4\right)\e^{-x}\le 0$
\end{enumerate}

\exo[Changement de variable]
On veut résoudre dans $\R$ l'équation $\e^{2x}+2\e^{x}-3=0$. On pose $X=\e^{x}$.
\begin{enumerate}
  \item Justifier que $\e^{2x}=X^{2}$.
  \item Résoudre l'équation $X^{2}+2X-3=0$.
  \item En déduire les solutions de l'équation de départ.
\end{enumerate}

\partiefiche{Étude de la fonction exponentielle}

\exo[Dériver une fonction qui contient $\e^{x}$]
Calculer la dérivée de chaque fonction ; factoriser par $\e^{x}$ quand c'est possible.
\begin{enumerate}[label=\alph*)]
  \item $f(x)=4\e^{x}+x^{2}$, sur $\R$
  \item $g(x)=(x+3)\,\e^{x}$, sur $\R$
  \item $h(x)=(3-x)\,\e^{x}$, sur $\R$
  \item $k(x)=\dfrac{\e^{x}}{x}$, sur $\intoo{0}{+\infty}$
  \item $m(x)=x^{2}\e^{x}$, sur $\R$
\end{enumerate}

\exo[Tangentes à la courbe de l'exponentielle]
On note $\mathcal{C}$ la courbe de la fonction exponentielle.
\begin{enumerate}
  \item Pour chaque nombre $a=-1$, $a=0$ et $a=2$, donner l'équation réduite de la tangente
        à $\mathcal{C}$ au point d'abscisse $a$.
  \item Pour chacune, calculer l'abscisse du point où elle coupe l'axe des abscisses.
  \item Écrire, par une phrase, la régularité observée.
\end{enumerate}

\exo[Courbe et tangente]
Pour tout nombre réel $x$, on pose $\varphi(x)=\e^{x}-x-1$.
\begin{enumerate}
  \item Calculer $\varphi'(x)$.
  \item Résoudre l'inéquation $\e^{x}-1\ge 0$ (écrire $1=\e^{0}$).
  \item Dresser le tableau de variations de $\varphi$ et donner son minimum.
  \item En déduire que, pour tout nombre réel $x$, $\e^{x}\ge x+1$.
  \item Interpréter ce résultat avec la courbe de l'exponentielle et sa tangente au
        point d'abscisse $0$.
\end{enumerate}

\exo[Comparer sans calculatrice]
\begin{enumerate}
  \item Pour chaque couple, ranger les deux nombres dans l'ordre croissant, en citant la
        propriété utilisée : $\e^{0{,}3}$ et $\e^{0{,}5}$ ; $\e^{-2}$ et $\e^{-3}$ ;
        $\e^{-0{,}1}$ et $1$.
  \item Pour $x\in\intff{-1}{2}$, encadrer $\e^{x}$.
  \item Quel est le signe de $\e^{x}-1$ selon les valeurs de $x$ ?
\end{enumerate}

\partiefiche{Dérivée de $\e^{u}$}

\exo[Dériver $\e^{u}$]
Calculer la dérivée de chaque fonction, définie sur $\R$.
\begin{enumerate}[label=\alph*)]
  \item $f(x)=\e^{3x}$
  \item $g(x)=\e^{-x}$
  \item $h(x)=\e^{2x-5}$
  \item $k(t)=50\,\e^{-0{,}2t}$
  \item $m(x)=\e^{x^{2}+1}$
  \item $p(x)=\e^{1-x^{2}}$
\end{enumerate}

\exo[Produit et exponentielle]
Calculer la dérivée de chaque fonction, définie sur $\R$, et la factoriser.
\begin{enumerate}[label=\alph*)]
  \item $f(x)=x\,\e^{3x}$
  \item $g(x)=(2x+1)\,\e^{-x}$
  \item $h(x)=(x+1)\,\e^{2x}$
  \item $k(x)=x^{2}\e^{-x}$
\end{enumerate}

\exo[Une dérivée proportionnelle à la fonction]
\begin{enumerate}
  \item Pour tout nombre réel $x$, on pose $f(x)=3\,\e^{-2x}$. Montrer que
        $f'(x)+2f(x)=0$.
  \item On pose $g(x)=\e^{4x+1}$. Montrer que $g'(x)=4g(x)$.
  \item Trouver une fonction $h$ telle que $h'(x)=5h(x)$ pour tout nombre réel $x$, et
        $h(0)=2$.
\end{enumerate}

\exo[Une courbe en cloche]
On considère la fonction $f$ définie sur $\intff{-3}{3}$ par $f(x)=\e^{-0{,}25x^{2}}$.
\begin{enumerate}
  \item Calculer $f'(x)$.
  \item Étudier le signe de $f'(x)$, puis dresser le tableau de variations de $f$.
  \item Quel est le maximum de $f$ ? Pour quelle valeur de $x$ est-il atteint ?
  \item Compléter un tableau de valeurs de $f$ sur $\intff{-3}{3}$ (au centième, pas de
        $1$), puis tracer la courbe de $f$ (unité : $2$~cm).
\end{enumerate}

\partiefiche{Étude de fonctions avec des exponentielles}

\exo[Étudier une fonction]
On considère la fonction $f$ définie sur $\intff{-2}{2}$ par $f(x)=(2x-3)\,\e^{x}$.
\begin{enumerate}
  \item Calculer $f'(x)$.
  \item Étudier le signe de $f'(x)$.
  \item Dresser le tableau de variations de $f$ et donner son minimum.
  \item Donner une équation de la tangente $T$ à $\Cf$ au point d'abscisse $0$.
  \item Compléter un tableau de valeurs (au dixième, pas de $1$), puis tracer $T$ et
        $\Cf$.
\end{enumerate}

\exo[Avec $\e^{-x}$]
On considère la fonction $f$ définie sur $\intff{-1}{4}$ par $f(x)=(x+1)\,\e^{-x}$.
\begin{enumerate}
  \item Montrer que $f'(x)=-x\,\e^{-x}$.
  \item Étudier le signe de $f'(x)$, puis dresser le tableau de variations de $f$.
  \item L'affirmation « pour tout $x$ de $\intff{-1}{4}$, $f(x)\le 1$ » est-elle vraie ?
        Justifier.
  \item Donner l'équation réduite de la tangente à $\Cf$ au point d'abscisse $1$.
\end{enumerate}

\exo[Un quotient]
On considère la fonction $f$ définie sur $\R$ par $f(x)=\dfrac{\e^{x}}{x^{2}+1}$.
\begin{enumerate}
  \item Montrer que $f'(x)=\dfrac{(x-1)^{2}\,\e^{x}}{\left(x^{2}+1\right)^{2}}$.
  \item En déduire le sens de variation de $f$.
  \item Donner une équation de la tangente à $\Cf$ au point d'abscisse $1$.
  \item Un élève affirme : « $f'(1)=0$, donc $f$ a un extremum en $1$. » Répondre par une
        phrase.
\end{enumerate}

\exo[Tangentes horizontales]
On considère la fonction $f$ définie sur $\R$ par $f(x)=\left(x^{2}-3\right)\e^{x}$.
\begin{enumerate}
  \item Résoudre l'équation $f(x)=0$.
  \item Montrer que $f'(x)=(x+3)(x-1)\,\e^{x}$.
  \item Donner les coordonnées des points de $\Cf$ où la tangente est horizontale.
  \item Dresser le tableau de variations de $f$ et donner ses extremums locaux.
\end{enumerate}

\partiefiche{Croissance et décroissance exponentielles}

\exo[Tracer $\e^{kt}$ et $\e^{-kt}$]
On considère les fonctions $f$ et $g$ définies sur $\intff{-4}{4}$ par
$f(t)=\e^{0{,}3t}$ et $g(t)=\e^{-0{,}3t}$.
\begin{enumerate}
  \item Donner le sens de variation de $f$ et celui de $g$.
  \item Compléter un tableau de valeurs de $f$ et de $g$ (au dixième, pas de $1$).
  \item Tracer les deux courbes dans un même repère (unité : $1$~cm).
  \item Laquelle modélise une croissance exponentielle ? une décroissance
        exponentielle ? Répondre par une phrase.
\end{enumerate}

\exo[Reconnaître une suite géométrique]
Pour chaque suite $(u_{n})$, $(v_{n})$, $(w_{n})$, montrer qu'elle est géométrique en
exprimant le terme suivant en fonction du terme courant ; donner sa raison, son premier
terme et son sens de variation.
\begin{enumerate}[label=\alph*)]
  \item $u_{n}=\e^{3n}$
  \item $v_{n}=4\,\e^{-0{,}5n}$
  \item $w_{n}=\e^{2n+1}$
\end{enumerate}

\exo[Élimination d'un médicament]
Un patient prend $80$~mg d'un médicament. La quantité restant dans son organisme,
en mg, $t$ heures plus tard, est $Q(t)=80\,\e^{-0{,}1t}$.
\begin{enumerate}
  \item Calculer $Q(0)$ et interpréter le résultat.
  \item Donner le sens de variation de $Q$.
  \item On note $u_{n}=Q(n)$. Montrer que $(u_{n})$ est géométrique ; donner sa raison,
        puis son arrondi au millième.
  \item De quel pourcentage la quantité baisse-t-elle chaque heure ?
  \item À l'aide d'un tableau de valeurs de la calculatrice, déterminer au bout de
        combien d'heures entières il reste moins de $20$~mg.
\end{enumerate}

\exo[Du modèle discret au modèle continu]
Une ville compte $2\,000$ habitants en 2026 ; sa population augmente de $4\,\%$ par an.
On note $u_{n}$ la population $n$ années après 2026.
\begin{enumerate}
  \item Justifier que $(u_{n})$ est géométrique ; exprimer $u_{n}$ en fonction de $n$.
  \item On cherche un modèle continu $f(t)=2\,000\,\e^{kt}$, avec $f(n)=u_{n}$ pour tout
        entier $n$. Justifier que $\e^{k}=1{,}04$.
  \item À l'aide d'un tableau de valeurs de la fonction exponentielle (pas de
        $0{,}001$), donner l'arrondi de $k$ au millième.
  \item Avec cette valeur, estimer la population au milieu de l'année 2028 ($t=2{,}5$),
        et la comparer à $u_{2}$ et $u_{3}$.
\end{enumerate}

\partiefiche{Problèmes concrets}

\exo[Bénéfice d'un atelier]
Un atelier fabrique entre $100$ et $2\,000$ trottinettes par mois. Pour $x$ centaines de
trottinettes, le bénéfice, en milliers d'euros, est
\[
  B(x)=(4x-20)\,\e^{-0{,}2x},\qquad x\in\intff{1}{20}.
\]
\begin{enumerate}
  \item Calculer $B(1)$, arrondi au dixième, et interpréter le résultat.
  \item Résoudre l'inéquation $B(x)\ge 0$ et interpréter.
  \item Montrer que $B'(x)=(8-0{,}8x)\,\e^{-0{,}2x}$.
  \item Dresser le tableau de variations de $B$.
  \item Combien de trottinettes faut-il fabriquer pour un bénéfice maximal ? Que vaut-il,
        à l'euro près ?
\end{enumerate}

\exo[Refroidissement d'un four à pain]
Un four à pain est éteint à 22~h ; sa température, en \textdegree C, $t$ heures plus tard,
est $T(t)=20+780\,\e^{-0{,}1t}$, pour $t\ge 0$.
\begin{enumerate}
  \item Quelle est la température du four quand on l'éteint ?
  \item Calculer $T'(t)$ ; en déduire le sens de variation de $T$.
  \item Calculer $T'(0)$ et l'interpréter (vitesse de refroidissement).
  \item On peut nettoyer le four quand sa température est inférieure à
        $100$~\textdegree C. Voici des valeurs de $T(t)$ arrondies au dixième.
        \begin{center}
        \renewcommand{\arraystretch}{1.3}
        \begin{tabular}{|c|c|c|c|c|c|}
        \hline
        $t$ & $20$ & $21$ & $22$ & $23$ & $24$\\ \hline
        $T(t)$ & $125{,}6$ & $115{,}5$ & $106{,}4$ & $98{,}2$ & $90{,}8$\\ \hline
        \end{tabular}
        \end{center}
        À partir de quelle heure entière peut-on le nettoyer ? Justifier.
\end{enumerate}

\exo[Le plus grand potager]
Un jardin est limité par les axes, la droite d'équation $x=5$ et la courbe $\Cf$ de
$f(x)=6\,\e^{-0{,}5x}$ (unité : le mètre). On y place un potager rectangulaire $OABC$,
avec $A$ d'abscisse $x$ sur l'axe des abscisses et $B$ sur $\Cf$.
\begin{center}
\begin{tikzpicture}[x=0.55cm,y=0.55cm]
  \repere{-1}{-1}{6}{7}
  \gradx{1,2,3,4,5} \grady{1,2,3,4,5,6}
  \draw[coulCourbe,line width=1.1pt,domain=0:5,samples=60]
    plot (\x,{6*exp(-0.5*\x)});
  \draw[coulExo,line width=0.9pt] (0,0) rectangle (1.5,2.834);
  \draw[dashed] (5,0)--(5,0.49);
  \node[font=\footnotesize,below] at (1.5,-0.3) {$A$};
  \node[font=\footnotesize,above right] at (1.5,2.834) {$B$};
  \node[font=\footnotesize,left] at (-0.3,2.834) {$C$};
  \node[coulCourbe,font=\footnotesize,right] at (3,1.6) {$\Cf$};
\end{tikzpicture}
\end{center}
\begin{enumerate}
  \item Justifier que l'aire du potager, en m$^{2}$, est $g(x)=6x\,\e^{-0{,}5x}$, pour
        $x\in\intff{0}{5}$.
  \item Montrer que $g'(x)=(6-3x)\,\e^{-0{,}5x}$.
  \item Dresser le tableau de variations de $g$.
  \item Où placer $A$ pour que le potager soit le plus grand ? Donner son aire, au
        dm$^{2}$ près.
\end{enumerate}

\end{multicols}

\end{document}
